\noindent Este trabalho investiga o impacto de trocas de elenco realizadas ao longo da temporada regular sobre o desempenho de equipes da \textit{National Basketball Association} (NBA), tratando a troca como um evento multidimensional, e não como algo redutível a um único indicador. Para isso, quatro hipóteses foram operacionalizadas e testadas. A primeira trata da associação entre a variação da desigualdade salarial em quadra — medida pelo Índice de Gini ponderado por minutos jogados, escolhido entre seis métricas concorrentes de desigualdade salarial — e o desempenho subsequente da equipe. A segunda avalia a associação entre o ajuste funcional do jogador adquirido, diagnosticado pela metodologia da Contribuição Relativa derivada do \textit{Win Probability Added} (WPA), e o impacto de curto prazo da troca. A terceira examina a associação entre a experiência dos jogadores adquiridos, medida pela idade média no momento da negociação, e esse mesmo impacto. A quarta investiga a associação entre a severidade da lacuna posicional preenchida e o tamanho do pacote de jogadores negociado. Os dados foram coletados por meio de uma arquitetura de dados em camadas (\textit{Medallion Lakehouse}), a partir de APIs públicas e fontes históricas da NBA, cobrindo treze temporadas, de 2012-13 a 2024-25. As hipóteses foram testadas por correlação de Spearman e pelo teste de Mann-Whitney \textit{U}, em amostras que variam de 282 a 370 casos conforme o teste. Os resultados encontram suporte estatístico, ainda que com efeito fraco, para três das quatro hipóteses. A desigualdade salarial associa-se negativamente ao desempenho subsequente ($\rho = -0{,}189$; $p = 0{,}0014$), favorecendo moderadamente a Teoria da Equidade sobre a Teoria do Torneio e concentrando-se nas equipes de perfil competitivo ofensivo. A experiência dos jogadores adquiridos associa-se positivamente ao impacto de curto prazo da troca ($\rho = +0{,}110$; $p = 0{,}0346$). O tamanho do pacote negociado, por sua vez, acompanha a severidade da lacuna preenchida ($\rho = -0{,}280$; $p < 0{,}0001$). Já o ajuste funcional do jogador adquirido não encontra suporte estatístico em sua forma categórica original, ainda que um sinal mais fraco sobreviva em sua versão contínua. Em conjunto, os quatro resultados sustentam a leitura de que a troca de elenco é, de fato, um evento multidimensional, no qual nenhuma dimensão isolada decide sozinha o sucesso da movimentação.

\vspace{0.5cm}
\noindent\textbf{Palavras-chave}: Trocas de elenco. Desigualdade salarial. Teoria da Equidade. Desempenho esportivo. NBA.
